\chapter{酸和碱不只是酸味和滑腻}

\begin{kaichang}
切一片柠檬，用舌尖碰一下——一股尖锐的酸味立刻让你眯起眼睛。再去做另一件事：洗手时故意多搓一会儿肥皂，注意手指之间那种滑溜溜、怎么冲都冲不干净的感觉。

现在请你猜两件事。第一：酸味是柠檬的“味道”，还是你的舌头在报告某种微观事件？第二：肥皂水的滑腻是肥皂本身很滑，还是它正在对你的皮肤做什么手脚？

还有一个更反常识的问题：把一块干冰（固态二氧化碳）放进水里，水会变酸；把氨水滴进水里，水会变“碱”——可干冰和氨气分子里，既没有盐酸那样的氢，也没有烧碱那样的氢氧根。它们凭什么？

这一章结束时，你会看到“酸”和“碱”这三个字背后，其实是一场永不停歇的质子接力赛。先把你对柠檬和肥皂的猜想写下来。
\end{kaichang}

\section{酸写在舌头上，碱写在皮肤上}

人类认识酸碱，是从身体开始的。醋、柠檬、酸梅、酸奶——凡是酸的，尝起来都酸。古人甚至给醋起了名字就叫“酸”（英语 acid 来自拉丁语“酸的”）。碱的认识稍微曲折一点：人们发现某些东西能洗掉油污、摸起来滑腻——草木灰泡的水、肥皂水、小苏打水。阿拉伯人把草木灰里的有效成分叫作 al-qaly，这就是英语 alkali（碱）的词根。

只靠舌头和手，人们积累了酸碱的第一批“档案”：

\begin{itemize}[itemsep=2pt]
\item 酸的溶液大多有酸味（\textbf{注意：这只是生活经验，绝对不能用嘴去尝实验室里的任何东西}，很多酸会烧伤口腔，很多酸的毒性跟酸味毫无关系）。
\item 酸能让某些有颜色的植物汁变色：紫甘蓝汁、牵牛花汁遇酸变红。
\item 酸滴在石灰石、贝壳、小苏打上会冒气泡（二氧化碳）。
\item 碱的溶液摸起来滑腻，能去油污，也能让紫甘蓝汁变绿甚至变黄。
\item 酸和碱混在一起，两边的“脾气”会互相抵消——醋倒进小苏打水里，嘶嘶冒完泡之后，既不酸也不滑了。
\end{itemize}

这些现象几百年里被反复使用：用酸洗金属、用草木灰制肥皂、用酸碱中和给胃“灭火”。但“档案”不等于理解。为什么所有的酸——盐酸、醋酸、柠檬酸、碳酸——尝起来都是同一种“酸”？它们的分子长得完全不同，凭什么给你同一种感觉？

要让不同的物质给出同一种感觉，最省事的解释是：它们在水里交出了\textbf{同一种粒子}。

\section{第一个定义：它们在水里放出什么}

1880 年代，瑞典一位年轻的化学博士生盯上了这个猜想。先记住结论，故事待会儿讲：他的名字叫阿伦尼乌斯，他提出——

\begin{qiaoliang}
\textbf{阿伦尼乌斯定义}：酸是在水中能放出氢离子 \chem{H^+} 的物质；碱是在水中能放出氢氧根离子 \chem{OH^-} 的物质。酸碱中和，就是 \chem{H^+} 和 \chem{OH^-} 结合成水。
\end{qiaoliang}

这个定义一下子解释了“为什么酸都是酸的”：盐酸、醋酸、柠檬酸在水里都会放出 \chem{H^+}，你的味蕾对酸的感知，本质上就是在探测 \chem{H^+} 的浓度。它同样解释了“为什么中和会互相抵消”：\chem{H^+} 和 \chem{OH^-} 一见面就抱成水分子，两边都“消失”了。

不过这里要立刻修正一个粒子层面的画面。水溶液里几乎不存在光着身子的 \chem{H^+}。氢原子失去唯一的电子后，剩下的是一个\textbf{裸质子}——一个又小又带电的粒子，电场强得吓人。在第 23 章我们讲过，水分子是极性的，氧端带负电；于是裸质子一出现，立刻被周围水分子的氧端“抱住”，形成水合氢离子 \chem{H_3O^+}（它还会拖着更多水分子一起行动）。所以，更诚实的写法是：

\[\chem{HCl} + \chem{H_2O} \rightarrow \chem{H_3O^+} + \chem{Cl^-}\]

\[\chem{NaOH} \rightarrow \chem{Na^+} + \chem{OH^-}\]

习惯上大家图省事写成 \chem{H^+}，你也经常会看到这样的简写。只要心里记着：酸溶液里真正巡游的“质子”，都是坐在水分子上的。

阿伦尼乌斯定义干净利落，但开场里的那个问题它答不上来：氨气 \chem{NH_3} 溶于水，明明是碱性的，能中和酸、能让指示剂变色——可氨分子里连一个氧原子都没有，上哪儿放 \chem{OH^-}？还有：把氯化氢气体和氨气两瓶气体瓶口相对，它们会直接反应生成一股白烟（氯化铵 \chem{NH_4Cl} 的微小颗粒）——整个过程没有水参与，按阿伦尼乌斯的说法，这连“酸碱反应”都算不上，可它看起来分明就是酸和碱在互相抵消。

定义太好用，又解释不了新事实，这正是理论该升级的信号。

\begin{zhengju}
阿伦尼乌斯是怎么想到“离解”的？靠的不是尝味道，而是两组当时说不通的实验数据。

第一组是导电。纯水的导电性极差，但往水里加盐、加酸、加碱，导电性立刻变好。当时的解释是“通电才把分子拆成了带电碎片”。可阿伦尼乌斯注意到一个别扭之处：溶液的导电性随浓度变化的方式，不像是“通电拆出来的”，倒像是带电粒子本来就在水里，通不通电都在。

第二组更硬。前面在第 25 章讲过，溶液的凝固点下降、渗透压大小，只数粒子的“个数”。实验发现，\ud{1}{mol} 食盐溶进水里，凝固点下降的幅度差不多是 \ud{1}{mol} 糖的两倍；氯化钙 \chem{CaCl_2} 更接近三倍。如果盐在水里还是“一整块一整块”的分子，凭什么数出两份、三份粒子？当时的权威化学家门捷列夫等人提出了替代解释：也许是溶质和水结合成了松散的“水合物”。争不出结果。

1884 年，阿伦尼乌斯在博士论文里提出一个大胆的统一答案：盐、酸、碱在水里\textbf{自己就拆成了带电的离子}，不需要通电。食盐变成 \chem{Na^+} 和 \chem{Cl^-} 两份粒子，所以凝固点按两份算；盐酸变成 \chem{H^+} 和 \chem{Cl^-}，所以能导电、有酸味。评审教授们觉得这个没名气的年轻人异想天开，论文只给了勉强及格的低分。

但阿伦尼乌斯不服气，他把论文寄给欧洲多位化学家，得到了范特霍夫、奥斯特瓦尔德等人的支持。关键的后援来自物理：1895 年汤姆孙发现电子，人们明白了“带电的原子”并不荒唐——原子可以失去或得到电子。到 1903 年，当年那篇勉强及格的论文为阿伦尼乌斯赢得了诺贝尔化学奖。这个故事值得记住的不是反转打脸，而是方法：他手里没有显微镜能看到离子，他靠的是\textbf{几个彼此独立的现象指向同一个解释}——导电、凝固点、酸味、中和，都要“溶液里预先存在离子”才讲得通。当好几条证据链都指向同一个看不见的粒子时，我们就敢给它发身份证。
\end{zhengju}

\section{更宽的定义：质子换手}

修补阿伦尼乌斯定义的工作，1923 年由丹麦人布朗斯特和英国人劳里各自独立完成。他们的想法是：别盯着“放出什么离子”，盯\textbf{质子本身的去向}。

\begin{qiaoliang}
\textbf{布朗斯特定义}：酸是\textbf{给出质子}（\chem{H^+}）的物质，碱是\textbf{接受质子}的物质。酸碱反应的本质是质子从一个粒子转移到另一个粒子，不要求在水里进行。
\end{qiaoliang}

用这个定义再看氨水，一切就顺了。氨分子自己不放 \chem{OH^-}，它从水分子那里\textbf{抢}一个质子：

\[\chem{NH_3} + \chem{H_2O} \rightleftharpoons \chem{NH_4^+} + \chem{OH^-}\]

氨是碱（接受质子），水在这里充当了酸（给出质子），被抢得只剩半个身子的水分子就变成了 \chem{OH^-}。注意箭头是可逆的——这是第 31 章讲过的动态平衡：绝大多数氨分子抢到质子后，铵根 \chem{NH_4^+} 又会把质子还回去，溶液里任意时刻只有一小部分是 \chem{NH_4^+} 和 \chem{OH^-}。这就是为什么氨水是“温和”的碱。

布朗斯特定义还顺手揭开了水的双面性格。遇到盐酸，水接受质子变成 \chem{H_3O^+}，水当碱；遇到氨，水给出质子变成 \chem{OH^-}，水当酸。水又酸又碱，取决于对手是谁。甚至水分子之间也会互相递质子：一个水分子把质子塞给邻居，自己变成 \chem{OH^-}，邻居变成 \chem{H_3O^+}——这叫水的自耦电离。纯水里每时每刻都有大约十亿分之几个水分子处于“拆了”的状态。这个数字小到可以忽略吗？恰恰相反，本章后半部分的 pH 全部建立在它头上。

还有一个更漂亮的推论。质子给出者交出质子后，剩下的部分随时可以把质子抢回来——所以每一种酸的身后都站着一个碱；碱拿到质子后也能再给出去——每一种碱的身后都站着一个酸。这样配成的一对，叫\textbf{共轭酸碱对}：醋酸的共轭碱是醋酸根，铵根的共轭碱是氨。酸越强，给出的质子越难要回来，它的共轭碱就越弱——像一场越打越弱的回音。

而开场提到的干冰，答案也浮出水面：\chem{CO_2} 本身不给质子，但它先和水反应生成碳酸 \chem{H_2CO_3}，碳酸再给质子。这条路径——\chem{CO_2} 溶进水里变酸——听着像化学课堂 trivia，其实正决定着海洋和血液的命运，本章后面会回来细讲。

最后，给看得更远的人一句话：还有第三个、最宽的定义——\textbf{路易斯定义}：酸是接受一对电子的粒子，碱是给出一对电子的粒子。它把布朗斯特定义整个装了进去（氨把孤对电子递给质子，就是路易斯碱），还能解释那些根本没有质子参与的反应，比如金属离子被一群水分子或氨分子团团围住形成配离子。那是第 35 章“配位”的主场，这里只立个路牌：三任定义，一任比一任宽，但不是后者推翻前者，而是像地图越画越大——在溶液里谈酸碱，布朗斯特的“质子换手”仍然是最称手的工具。

\section{强和浓，是两回事}

接下来是两个被混用了一辈子的词。请先记住红线：\textbf{“强弱”说的是性格，“浓稀”说的是数量，这是两组完全独立的概念。}

酸有强弱，说的是它在水里交质子交得\textbf{彻不彻底}。盐酸是强酸：\chem{HCl} 分子进了水，几乎全部拆成 \chem{H_3O^+} 和 \chem{Cl^-}，一个不留。醋酸是弱酸：\chem{CH_3COOH} 分了水之后，只有百分之一左右的分子交出质子，其余百分之九十九保持原样，反应是个停在半路的平衡：

\[\chem{CH_3COOH} + \chem{H_2O} \rightleftharpoons \chem{CH_3COO^-} + \chem{H_3O^+}\]

浓稀说的完全是另一件事：一升水里你\textbf{放了多少}酸。\ud{10}{mol} 醋酸溶进一升水，是浓醋酸；\ud{0.0001}{mol} 盐酸溶进一升水，是稀盐酸。于是四种组合都存在，而且生活中全都能见到：

\begin{table}[htbp]
\centering
\begin{tabular}{lll}
\toprule
 & 浓（溶质多） & 稀（溶质少） \\
\midrule
强（电离彻底） & 浓硫酸、浓盐酸（实验室试剂瓶） & 胃里的盐酸（约 \ud{0.03}{mol\,L^{-1}}） \\
弱（电离不彻底） & 冰醋酸（食醋的浓缩老祖宗） & 白醋（约 5\% 的醋酸水溶液） \\
\bottomrule
\end{tabular}
\end{table}

混淆这两组概念会出真危险。有人听说醋酸是“弱酸”，就以为冰醋酸很安全——错了，冰醋酸是浓的弱酸，照样能烧伤皮肤；“弱”只说明它的质子交得不痛快，可架不住它量大。反过来，极稀的盐酸是强酸，却已经淡到可以存在于你的胃里。弱不等于温和，强不等于危险，危险的是\textbf{具体浓度下的总量}——这和第 1 章“剂量决定毒性”是同一条原则。

顺便提醒一条实验室铁律，以后遇到别忘：稀释浓硫酸时，永远把酸缓缓倒入水中并搅拌，绝不能反过来。浓硫酸遇水放出大量的热（第 27 章的能量账），水倒进酸里会在表面沸腾飞溅，把滚烫的酸甩到你脸上。

\begin{wuqu}
“pH 小于 7 的东西会腐蚀，碰不得；碱性东西温和，滑滑的很安全。”——两个半句都错。前半句的反例：你的胃液 pH 约 1.5，可乐 pH 约 2.5，都在你体内或被你喝下去。后半句更危险：强碱的腐蚀性常常比强酸更阴险。酸烧伤立刻刺痛，你会马上躲开；而强碱（比如疏通下水道的烧碱溶液）破坏皮肤时\textbf{不太疼}——那种滑腻感，其实是碱在把你皮肤表层的油脂分解成肥皂（这个过程就叫“皂化”），等你察觉不对，伤害已经做下了。所以通渠剂、烤箱清洁剂这类强碱用品，务必戴手套，而且绝对不能用滑腻感去“鉴别”任何化学品。判断酸碱性靠指示剂和 pH 计，不靠皮肤和舌头。
\end{wuqu}

\section{用对数给质子记账}

现在来给酸度定量。溶液里 \chem{H^+}（记住，其实是 \chem{H_3O^+}）的浓度范围大得吓人：浓盐酸里可达每升约 \ud{10}{mol}，浓烧碱溶液里低到每升 $10^{-15}$ mol——跨越 16 个数量级。在横跨万亿倍的尺度上记账，最顺手的工具是对数：不看数字本身，只看 10 的指数。

\begin{qiaoliang}
\[\mathrm{pH} = -\lg\,[\chem{H^+}]\]
其中 $[\chem{H^+}]$ 是每升溶液中氢离子的摩尔数（严格说是“活度”，见本章挑战层）。pH 每差 1，氢离子浓度差 10 倍；每差 2，差 100 倍。浓度越高，pH 反而越低——因为前面有个负号。
\end{qiaoliang}

拿尺子量一量身边的世界：柠檬汁 pH 约 2.4，可乐约 2.5，白醋约 2.9，牛奶约 6.5，纯水（$25\,^{\circ}\mathrm{C}$）7.0，血液 7.4，肥皂水约 10，通渠剂可达 13 以上。

来做一个数量级估算，感受一下对数尺度的夸张。胃酸的 pH 约 1.5，血液的 pH 约 7.4，相差 5.9 个单位。氢离子浓度差 $10^{5.9}$ 倍，约等于 80 万倍——\textbf{你的胃里每升液体中的氢离子，是血液里的将近一百万倍}。这两种液体同处一个身体，只隔着一层胃壁。胃壁凭什么不被自己消化？靠黏液和快速更新的上皮细胞筑起防线；防线一旦失守，就是胃溃疡。再看另一个方向：海水 pH 从工业革命前的约 8.2 降到今天的约 8.1，只降了 0.1——但换算回来，氢离子浓度升高了约 $26\%$（$10^{0.1}\approx1.26$）。在对数尺度上，0.1 不是小数目。这 0.1 的账，后面要算到海洋头上。

\section{pH 7 不是宇宙的常数}

几乎每个初中生都背过：“pH 等于 7 是中性。”请把这句话改掉——\textbf{它只在 $25\,^{\circ}\mathrm{C}$ 附近成立，不是普适定律}。

回到中性的真正定义：中性不是“pH 等于某个数”，而是 $[\chem{H^+}]=[\chem{OH^-}]$，氢离子和氢氧根一样多。纯水永远自带这个平衡：水分子互相递质子，生成一对 \chem{H_3O^+} 和 \chem{OH^-}，两者浓度相等——所以\textbf{任何温度下的纯水都是中性的}，这一点没有例外。

但“相等”的那个数值，随温度漂移。在第 31 章我们学过：平衡会朝着抵消外界改变的方向移动。水分子的自耦电离是吸热过程（把一个水分子拆成两个离子要付能量），于是升温时平衡向电离方向挪，\chem{H^+} 和 \chem{OH^-} 同时变多。定量地说：$25\,^{\circ}\mathrm{C}$ 时纯水的 $[\chem{H^+}]=10^{-7}$ mol/L，恰好 pH 7.0；升温到 $100\,^{\circ}\mathrm{C}$，$[\chem{H^+}]$ 涨到约 $10^{-6.1}$ mol/L，中性点变成 pH 约 6.1；降温到 $0\,^{\circ}\mathrm{C}$，中性点退到 pH 约 7.5。

所以一杯 $100\,^{\circ}\mathrm{C}$ 的纯水，pH 是 6.1，它依然中性——因为 \chem{H^+} 和 \chem{OH^-} 还是一样多。反过来，如果在 $100\,^{\circ}\mathrm{C}$ 时把溶液调到 pH 恰好 7.0，它其实是\textbf{碱性}的：氢离子比这个温度下的中性值少，说明氢氧根更多。背“7 等于中性”，等于把温度计上的一个刻度当成物理定律。真正不变的规律是：中性由两种离子的相对多少定义，而 pH 只是 \chem{H^+} 这一边的读数。

\begin{shiyan}
\textbf{紫甘蓝：你家厨房里的 pH 计}

材料：紫甘蓝（紫包菜）几片叶子、热水、五个透明杯子、白醋、小苏打水、肥皂水、雪碧（或其他无色碳酸饮料）、纯净水。请一位成年人协助倒热水。

步骤：
\begin{enumerate}[itemsep=1pt]
\item 紫甘蓝叶撕碎，用热水浸泡十分钟，得到深紫色的汁，放凉。紫色来自一类叫花青素的分子，它的结构会随抓住几个质子而改变，颜色随之改变。
\item 五个杯子各倒一点甘蓝汁，分别加入少量白醋、小苏打水、肥皂水、雪碧、纯净水。
\item 记录颜色，按“红—紫—蓝—绿—黄”的顺序排队。
\end{enumerate}

观察点：白醋杯应偏红，雪碧杯偏粉红，纯净水保持紫色，小苏打水偏蓝绿，肥皂水偏绿。同一份指示剂给出五种颜色，说明这五种液体的 \chem{H^+} 浓度按数量级排列——你刚刚用肉眼做了一次对数测量。

安全提示：本实验所有材料都是食品或日用品，安全；但\textbf{实验后的液体一律倒掉不喝}（实验器皿不是餐具），热水防烫。想升级的话，可以试试牵牛花汁、紫薯水、红茶（加柠檬会变色）。
\end{shiyan}

\section{中和、缓冲与滴定：三样称手的工具}

酸碱知识落到地面上，主要靠三件工具。

\textbf{中和}最直白：\chem{H_3O^+} 和 \chem{OH^-} 相遇生成水，同时放热——把等量的稀盐酸和稀烧碱溶液混合，温度会明显升高，这是第 27 章能量账里经典的一笔：

\[\chem{H_3O^+} + \chem{OH^-} \rightarrow \chem{2H_2O}\]

胃酸过多时吃的抗酸药就是中和的应用：碳酸氢钠、碳酸钙、氢氧化铝 \chem{Al(OH)_3} 各显神通，用碱性把多余的 \chem{H^+} 按回去。注意剂量原则：偶尔烧心用一点无妨，长期依赖抗酸药会掩盖真正的病因，还可能打乱胃的酸碱节律——这类问题交给医生，别自己当自己的药剂师。

\textbf{缓冲液}更精妙。中和是“来多少灭多少”，缓冲是“不管来多少，都尽量不动声色”。以你血液里的碳酸—碳酸氢根缓冲对为例，它是一对共轭酸碱联手值班：

\[\chem{H_2CO_3} + \chem{H_2O} \rightleftharpoons \chem{HCO_3^-} + \chem{H_3O^+}\]

外来酸闯进血液，\chem{HCO_3^-} 把它接住变成碳酸；外来碱闯进来，碳酸给出质子把它中和。按第 31 章的勒夏特列原理，这个平衡左右腾挪，把血液的 pH 死死按在 7.35～7.45 之间。你的细胞只能在这个窄到 0.1 的窗口里正常工作——酶的形状、离子的电荷全都对酸碱度敏感。血液 pH 跌到 7.2 以下或升到 7.6 以上会危及生命。顺便说，碳酸还有一条出路：分解成 \chem{CO_2} 从肺里呼出去。所以呼吸深浅也参与调节酸碱——憋气一分钟，血液会轻微变酸，因为你囤住了 \chem{CO_2}。呼吸系统和缓冲系统是一套组合拳，这个伏笔我们先埋在这里。

\textbf{滴定}是化学侦探的定量手艺。想知道一瓶醋里到底有多少醋酸？往醋里滴加浓度已知的烧碱溶液，同时用指示剂（或 pH 计）监视。酚酞在酸里无色、pH 超过约 8.2 变粉红——粉红刚冒头的那一刻，说明酸恰好被一滴不多一滴不少地中和完。记下用掉的碱液体积，按反应式 1:1 的比例一算，醋酸的总量就出来了。滴定的精髓是\textbf{用精确的体积和已知的浓度，反推未知的量}，这套“按反应比例记账”的思路，到第 35 章还会扩展成更全面的化学侦探术。

\section{牙齿、海洋和血液：三笔没算完的账}

把 pH 的眼光带出教室，有三笔账正在你身边发生。

第一笔在你的嘴里。牙釉质的主要成分是羟基磷灰石——一种含钙、磷酸根和氢氧根的矿物。酸能溶解它：当口腔 pH 降到约 5.5 以下，釉质开始脱矿。口腔里的细菌把糖发酵成酸，所以真正伤牙的不是糖本身，而是糖喂出来的酸；可乐更惨，它自带 pH 2.5 还附赠糖。好在唾液本身是缓冲液，能把 pH 拉回安全线，氟化物还能帮釉质重新矿化。频繁地小口慢喝碳酸饮料，等于让牙齿整天泡在酸浴里——伤害取决于\textbf{接触的总剂量}，又一次回到剂量原则。

第二笔在整个海洋。大气里的 \chem{CO_2} 不断溶进海水，走的就是本章开场干冰那条路：先生成碳酸，再给出质子。工业革命以来，人类烧掉的化石燃料让大气 \chem{CO_2} 多了约一半，海水表层 pH 相应从约 8.2 滑到 8.1——前面算过，这是氢离子浓度上涨约 $26\%$。麻烦在于，珊瑚和贝类用碳酸钙 \chem{CaCO_3} 盖房子，而多出来的 \chem{H^+} 会抢走碳酸根 \chem{CO_3^{2-}}（变成 \chem{HCO_3^-}），等于把建筑工地的砖先抽走了一批。这就是“海洋酸化”：海洋不会真的变成“酸”（pH 仍在 8 以上），但只要朝酸的方向滑，靠碳酸钙生存的生命就要多付成本。

第三笔在你的血管里，而且永不结账。你每跑一步，肌肉都在产生乳酸和 \chem{CO_2}；血液缓冲对、肺、肾三班倒，把 pH 钉在 7.4 附近。这套系统调的是第 31 章的化学平衡，守的是第 32 章的酸碱，到了第二册讲代谢时，你还会看到它怎样被写进生命的操作系统。

\section{回到柠檬和肥皂}

现在兑现开场的承诺。

柠檬的酸味，是果肉细胞里储存的柠檬酸在水中给出质子、\chem{H_3O^+} 撞上你舌头上感受器的结果。你的舌头不是在“品尝柠檬”，而是在做 pH 测量——这套传感器古老到所有脊椎动物都在用，因为“体液酸碱度”对一切生命都是生死攸关的读数。顺带解开一个小谜团：为什么柠檬酸是一种弱酸，柠檬却酸得这么冲？因为它\textbf{浓}——柠檬汁的 pH 低到 2.4，靠的是量大，又一次印证强弱与浓稀无关。

肥皂的滑腻则是双重戏法。肥皂本身是油脂加碱皂化出来的脂肪酸盐，它的分子一头亲水一头亲油（第 22 章的分子间作用力在这里重逢），能拆散油污；而你摸肥皂水时觉得皮肤变滑，部分原因是残余的碱性真的在把你皮肤表面的油脂现场再皂化一遍——你摸到的滑，有一小部分是你自己的油脂变成的肥皂。这就是为什么用完强碱性的清洁剂，手会干涩发紧：皮脂被洗劫一空。滑腻不是肥皂的性格，是碱在你皮肤上开工的声音。

\section{质子会换手，电子也会}

回头看这一章，我们其实只讲了一种动作：一个带正电的粒子（质子）从一个分子手里转到另一个分子手里。酸是给出者，碱是接受者，中和是物归原主，缓冲是随时待命的接力队。

但化学里会“换手”的不止质子。还有另一种更贵重、更危险的粒子也在分子之间不停转移——电子。铁生锈，是铁把电子转给了氧；电池放电，是电子被迫绕路走了一圈外电路；你呼吸、吃饭、烧掉一块糖，归根到底都是电子从一种物质流向另一种物质。质子的转移定义了酸碱，电子的转移将定义下一章的主角：氧化还原。

\begin{tiaozhan}
（跳过不影响主线）pH 定义里的方括号，严格说不是浓度而是\textbf{活度}——粒子在拥挤溶液里的“有效浓度”。在稀溶液里活度约等于浓度，所以中学阶段直接用浓度不会出错；但在浓溶液或海水这种“离子汤”里，离子互相牵制，活度明显小于浓度，这正是精确测海水 pH 要用特殊电极和校正曲线的原因之一。

再来推一遍中性点随温度的漂移。定义水的离子积 $K_w=[\chem{H^+}][\chem{OH^-}]$。实验测得：$0\,^{\circ}\mathrm{C}$ 时 $K_w\approx1.1\times10^{-15}$，$25\,^{\circ}\mathrm{C}$ 时 $K_w=1.0\times10^{-14}$，$100\,^{\circ}\mathrm{C}$ 时 $K_w\approx5.1\times10^{-13}$。纯水中两者相等，所以 $[\chem{H^+}]=\sqrt{K_w}$。代入 $100\,^{\circ}\mathrm{C}$：$[\chem{H^+}]\approx7.1\times10^{-7}$ mol/L，pH $=7-\lg 7.1\approx6.14$；代入 $0\,^{\circ}\mathrm{C}$：pH $\approx7.48$。三个温度，三个“中性 pH”——而水在三个温度下都严格中性。$K_w$ 随温度上升，正是因为电离吸热、平衡正向移动；这不是背出来的表格，是勒夏特列原理的一次定量胜利。
\end{tiaozhan}

\section*{练一练}

\begin{sikaoti}
\item 画粒子图：各画一个烧杯，分别表示浓盐酸和浓醋酸（假设两份溶液里溶质分子总数相同）。用圆圈表示未电离的分子，用成对的带电荷小球表示离子，画出两杯溶液里粒子种类的差别。图旁注明：图中粒子大小和数量比例是示意，不代表真实比例。
\item 有人把 \ud{0.1}{mol} 盐酸和 \ud{0.1}{mol} 醋酸分别溶进 \ud{1}{L} 水里。两份溶液的浓度谁大？酸性强弱谁大？和足量小苏打反应，最终放出的 \chem{CO_2} 谁多？（提示：最后一问想想勒夏特列原理——弱酸不断被消耗，平衡会怎么走。）
\item 雨水的“正常”pH 约 5.6，并不是 7。用本章和上一节的内容解释：即使完全没有污染，雨水为什么也是弱酸性的？pH 5.6 的雨水，氢离子浓度是纯水（$25\,^{\circ}\mathrm{C}$）的多少倍？
\item 一杯纯水从 $25\,^{\circ}\mathrm{C}$ 加热到 $80\,^{\circ}\mathrm{C}$。它的 pH 会上升、下降还是不变？它变成酸性了吗？用 $[\chem{H^+}]$ 与 $[\chem{OH^-}]$ 的关系说明你的答案，并指出“pH 7 为中性”这个说法在哪个环节出了错。
\item 画物质流图：某人剧烈呕吐，丢失大量胃酸（盐酸）。用箭头和共轭酸碱对的知识，预测他的血液 pH 会向哪个方向偏移，血液的碳酸缓冲对会往哪边移动来“救火”？（提示：这是真实的临床问题，叫代谢性碱中毒；图里至少出现 \chem{H_2CO_3}、\chem{HCO_3^-} 和 \chem{CO_2} 三个角色。）
\end{sikaoti}
